\section{组织与落地：从原型团队到 battle-hardened 团队}

讲者在后半段强调了组织问题：企业内部出现了新的角色组合，包括 AI architect、agent designer、运营专家、语音体验角色（如 voice sommelier）和产品工程协同岗位。这些角色不是锦上添花，而是把 Agent 从“项目”变成“能力平台”的必要条件。

\subsection{职责重构与协作接口}

传统客服组织按渠道分工，AI 团队按模型能力分工。落地阶段需要新增“策略与流程产品经理”“对话质量运营”“安全与评测工程师”等横向角色，确保业务目标与技术目标在同一发布节奏上推进。

\begin{knowledgebox}{建议的最小跨职能单元}
\begin{itemize}[leftmargin=1.5em]
    \item 业务 owner：定义高价值问题域与成功标准；
    \item Agent engineer：编排工具、策略、记忆与回退；
    \item Eval engineer：维护 simulation、回归集、上线门槛；
    \item Safety owner：维护 guardrail、红队与审计策略；
    \item Ops analyst：追踪升级原因并驱动知识回灌。
\end{itemize}
\end{knowledgebox}

\subsection{一个可执行的 90 天推进框架}

\begin{table}[H]
\centering
\caption{面向 customer-facing agent 的 90 天落地节奏建议}
\begin{tabularx}{\textwidth}{>{\raggedright\arraybackslash}p{2.2cm} >{\raggedright\arraybackslash}X}
\toprule
阶段 & 关键任务 \\
\midrule
0--30 天 & 明确高价值场景与红线策略；建立最小工具链（身份校验、订单查询、状态更新）；搭建离线评测与对话回放 \\
31--60 天 & 接入 memory 与 warm start；上线 tracing 与升级归因；引入多轮 simulation，建立发布门槛 \\
61--90 天 & 渠道统一与人工接管优化；部署分层 guardrails；将 expert answers 回灌流程制度化，进入持续迭代 \\
\bottomrule
\end{tabularx}
\end{table}

\begin{importantbox}{课程级结论}
Agent 的竞争力不在“谁先接入最新模型”，而在“谁先建立稳定的迭代操作系统”。组织设计、评测纪律和安全治理决定了长期上限。
\end{importantbox}

\subsection{本章小结}

customer-facing agent 是跨职能系统工程。团队若只补模型工程能力，不补评测、安全、运营与产品协同能力，项目会在真实业务压力下停滞。

